\documentclass[11pt]{article}
\usepackage[margin=1in]{geometry}
\usepackage{amsmath,amssymb,mathtools,bm}

\newcommand{\Exp}{\operatorname{Exp}}
\newcommand{\SO}{\mathrm{SO}}
\newcommand{\Ad}{\mathbf{Ad}}
\newcommand{\R}{\mathbb{R}}

\title{Exercise 1: Jacobian of $R^{-1}p$}
\author{}
\date{}

\begin{document}
\section*{Problem}
Let
\[
R\in\SO(3),\qquad p\in\R^3.
\]
Calculate
\[
\boxed{\frac{\partial(R^{-1}p)}{\partial R}}
\]
using both right and left perturbation models.
\section{Right perturbation}

Use
\[
R'=R\Exp(\delta\theta).
\]

\subsection{Method 1: Chain rule}

Consider
\[
R\longmapsto R^{-1}\longmapsto R^{-1}p.
\]
Then
\[
{}^{R}J_R^{R^{-1}p}
=
{}^{R^{-1}}J_{R^{-1}}^{R^{-1}p}{}^{R}J_R^{R^{-1}}.
\]

For $S\in\SO(3)$, the right-perturbation Jacobian of $Sp$ is
\[
{}^{S}J_S^{Sp}=-Sp^\wedge.
\]
Setting $S=R^{-1}$,
\[
{}^{R^{-1}}J_{R^{-1}}^{R^{-1}p}=-R^{-1}p^\wedge.
\]

For inversion under a right perturbation,
\[
{}^{R}J_R^{R^{-1}}=-\Ad_R=-R.
\]

Therefore
\[
{}^{R}J_R^{R^{-1}p}
=(-R^{-1}p^\wedge)(-R)
=R^{-1}p^\wedge R.
\]

Using
\[
(Ra)^\wedge=Ra^\wedge R^T,
\]
with $R^{-1}$ in place of $R$,
\[
R^{-1}p^\wedge R=(R^{-1}p)^\wedge.
\]

Hence
\[
\boxed{
{}^{R}J_R^{R^{-1}p}=(R^{-1}p)^\wedge
\qquad\text{(right perturbation)}.
}
\]

\subsection{Method 2: Directly from the definition}

Choose an arbitrary $u\in\R^3$ and let the perturbation be $hu$.
Then
\[
{}^{R}J_R^{R^{-1}p}u
=
\lim_{h\to0}
\frac{
(R\Exp(hu))^{-1}p-R^{-1}p
}{h}.
\]

Since
\[
(R\Exp(hu))^{-1}
=
\Exp(-hu)R^{-1},
\]
we obtain
\[
{}^{R}J_R^{R^{-1}p}u
=
\lim_{h\to0}
\frac{
\Exp(-hu)R^{-1}p-R^{-1}p
}{h}.
\]

Using
\[
\Exp(-hu)=I-hu^\wedge+O(h^2),
\]
\[
{}^{R}J_R^{R^{-1}p}u
=
\lim_{h\to0}
\frac{
-hu^\wedge R^{-1}p+O(h^2)
}{h}.
\]

Since
\[
a^\wedge b=-b^\wedge a,
\]
\[
-u^\wedge R^{-1}p=(R^{-1}p)^\wedge u.
\]

Thus
\[
{}^{R}J_R^{R^{-1}p}u=(R^{-1}p)^\wedge u.
\]

Because $u$ is arbitrary,
\[
\boxed{
{}^{R}J_R^{R^{-1}p}=(R^{-1}p)^\wedge.
}
\]
\section{Left perturbation}

Use the left-plus convention
\[
R'=\Exp(\delta\theta)R.
\]

\subsection{Method 1: From equation (57)}

The left-perturbation result can also be obtained directly from the
right-perturbation Jacobian by using equation (57) of the paper:
\[
\boxed{
{}^{\varepsilon}\frac{Df(X)}{DX}
=
\mathbf{Ad}_{f(X)}
\;
{}^{X}\frac{Df(X)}{DX}
\;
\mathbf{Ad}_{X}^{-1}.
}
\tag{57}
\]

For the present exercise,
\[
f(R)=R^{-1}p.
\]

The right-perturbation Jacobian has already been found to be
\[
\boxed{
{}^{R}J_R^{R^{-1}p}
=
(R^{-1}p)^\wedge.
}
\tag{11}
\]

Applying (57) gives
\[
{}^{\varepsilon}J_R^{R^{-1}p}
=
\mathbf{Ad}_{R^{-1}p}
\;
{}^{R}J_R^{R^{-1}p}
\;
\mathbf{Ad}_{R}^{-1}.
\tag{12}
\]

The output
\[
R^{-1}p\in\mathbb{R}^3
\]
belongs to the additive translation group $(\mathbb{R}^3,+)$.  This
group is commutative, and its adjoint matrix is the identity:
\[
\boxed{
\mathbf{Ad}_{R^{-1}p}=I.
}
\tag{13}
\]

For $\SO(3)$,
\[
\boxed{
\mathbf{Ad}_{R}=R,
\qquad
\mathbf{Ad}_{R}^{-1}=R^{-1}.
}
\tag{14}
\]

Substituting (11)--(14) into (12),
\begin{align}
{}^{\varepsilon}J_R^{R^{-1}p}
&=
I\,(R^{-1}p)^\wedge R^{-1}
\\
&=
(R^{-1}p)^\wedge R^{-1}.
\end{align}

Using
\[
(R^{-1}p)^\wedge
=
R^{-1}p^\wedge R,
\]
we obtain
\begin{align}
{}^{\varepsilon}J_R^{R^{-1}p}
&=
R^{-1}p^\wedge RR^{-1}
\\
&=
\boxed{
R^{-1}p^\wedge.
}
\end{align}

Thus equation (57) converts the right-perturbation result directly into
the left-perturbation result:
\[
\boxed{
{}^{R}J_R^{R^{-1}p}
=
(R^{-1}p)^\wedge
\quad\Longrightarrow\quad
{}^{\varepsilon}J_R^{R^{-1}p}
=
R^{-1}p^\wedge.
}
\]

This agrees with the result obtained independently from the direct
left-perturbation limit.
\subsection{Method 2: Directly from the definition}


To make the Jacobian limit rigorous, choose an arbitrary direction
\[
u\in\R^3
\]
and set
\[
\delta\theta=hu.
\]

Then
\[
{}^{\varepsilon}J_R^{R^{-1}p}u
=
\lim_{h\to0}
\frac{
(\Exp(hu)R)^{-1}p-R^{-1}p
}{h}.
\]

Using
\[
(AB)^{-1}=B^{-1}A^{-1},
\]
we have
\[
(\Exp(hu)R)^{-1}
=
R^{-1}\Exp(-hu).
\]

Therefore
\[
{}^{\varepsilon}J_R^{R^{-1}p}u
=
\lim_{h\to0}
\frac{
R^{-1}\Exp(-hu)p-R^{-1}p
}{h}.
\]

For small $h$,
\[
\Exp(-hu)=I-hu^\wedge+O(h^2).
\]

Hence
\begin{align}
{}^{\varepsilon}J_R^{R^{-1}p}u
&=
\lim_{h\to0}
\frac{
R^{-1}(I-hu^\wedge)p-R^{-1}p+O(h^2)
}{h}
\\
&=
\lim_{h\to0}
\frac{
-hR^{-1}u^\wedge p+O(h^2)
}{h}.
\end{align}

Using
\[
u^\wedge p=-p^\wedge u,
\]
we get
\[
-R^{-1}u^\wedge p
=
R^{-1}p^\wedge u.
\]

Therefore
\[
{}^{\varepsilon}J_R^{R^{-1}p}u
=
R^{-1}p^\wedge u.
\]

Since $u$ is arbitrary,
\[
\boxed{
{}^{\varepsilon}J_R^{R^{-1}p}
=
R^{-1}p^\wedge.
}
\]
\section{Final comparison}

For the right perturbation
\[
R'=R\Exp(\delta\theta_R),
\]
\[
\boxed{
\frac{\partial(R^{-1}p)}{\partial R}
=
(R^{-1}p)^\wedge.
}
\]

For the left perturbation
\[
R'=\Exp(\delta\theta_L)R,
\]
\[
\boxed{
\frac{\partial(R^{-1}p)}{\partial R}
=
R^{-1}p^\wedge.
}
\]

These two Jacobians are related by
\[
(R^{-1}p)^\wedge
=
R^{-1}p^\wedge R,
\]
and the perturbation coordinates satisfy
\[
\delta\theta_L=R\,\delta\theta_R.
\]

Indeed,
\[
R^{-1}p^\wedge\delta\theta_L
=
R^{-1}p^\wedge R\,\delta\theta_R
=
(R^{-1}p)^\wedge\delta\theta_R.
\]

Thus the two perturbation conventions describe the same first-order
variation using different tangent-coordinate representations.
\section*{Useful identity}

For $R\in\SO(3)$ and $a\in\R^3$,
\[
\boxed{(Ra)^\wedge=Ra^\wedge R^T.}
\]

\end{document}
